\documentclass[11pt,a4paper]{article}
\usepackage[utf8]{inputenc}
\usepackage[T1]{fontenc}
\usepackage[spanish]{babel}
\usepackage{amsmath, amssymb}
\usepackage{booktabs}
\usepackage{geometry}
\geometry{margin=2.5cm}

\begin{document}

Entendido perfectamente. Tienes toda la razón: definirlo como «construir el siguiente evento» fue una simplificación imprecisa.

En un espacio muestral formal $(\Omega, \mathcal{F}, P)$, donde cada combinación elemental $x \in \Omega$ tiene una probabilidad clásica laplaciana uniforme:

$$P(x) = \frac{1}{|\Omega|}$$

la probabilidad de cualquier sexteto específico es fija e idéntica. Para clarificar qué está haciendo realmente el algoritmo matemáticamente, \textbf{reduciremos el espacio muestral a su mínima expresión}.

\medskip
\hrule
\medskip

\subsection*{1. El Espacio Muestral Reducido (Micro-Lotería)}

Imaginemos una lotería en miniatura:
\begin{itemize}
    \item Se eligen \textbf{2 números} de un total de \textbf{4 disponibles}: $\{1, 2, 3, 4\}$.
    \item Sin orden y sin reemplazo.
\end{itemize}

El espacio muestral $\Omega$ está formado por todas las combinaciones posibles:

$$\Omega = \big\{ \{1,2\}, \{1,3\}, \{1,4\}, \{2,3\}, \{2,4\}, \{3,4\} \big\}$$

La cardinalidad total es:

$$|\Omega| = \binom{4}{2} = 6$$

Dado que el mecanismo es uniforme, para todo evento elemental $x \in \Omega$:

$$P(x) = \frac{1}{6} \approx 16.67\%$$

Ningún par de números es más probable que otro a nivel microstate. $\{1,2\}$ tiene exactamente la misma probabilidad de salir que $\{3,4\}$.

\medskip
\hrule
\medskip

\subsection*{2. Variables Aleatorias y Partición del Espacio Muestral}

Aquí es donde entra la matemática detrás del algoritmo. Definamos una \textbf{variable aleatoria} $S(x)$ que representa la \textbf{suma de los números}:

$$S(x) = r_1 + r_2$$

Evaluemos el valor de $S$ para cada elemento de $\Omega$:

\begin{center}
\begin{tabular}{lc}
\toprule
\textbf{Evento elemental $x \in \Omega$} & \textbf{Suma $S(x)$} \\
\midrule
$\{1, 2\}$ & $3$ \\
$\{1, 3\}$ & $4$ \\
$\{1, 4\}$ & $5$ \\
$\{2, 3\}$ & $5$ \\
$\{2, 4\}$ & $6$ \\
$\{3, 4\}$ & $7$ \\
\bottomrule
\end{tabular}
\end{center}

Ahora calculemos la probabilidad no de un boleto específico, sino de los \textbf{subconjuntos inducidos por la suma} ($E_s = \{x \in \Omega \mid S(x) = s\}$):

\begin{itemize}
    \item Para suma 3: $|E_3| = 1 \implies P(S=3) = \frac{1}{6} \approx 16.67\%$
    \item Para suma 4: $|E_4| = 1 \implies P(S=4) = \frac{1}{6} \approx 16.67\%$
    \item \textbf{Para suma 5: $|E_5| = 2 \implies P(S=5) = \frac{2}{6} \approx 33.33\%$}
    \item Para suma 6: $|E_6| = 1 \implies P(S=6) = \frac{1}{6} \approx 16.67\%$
    \item Para suma 7: $|E_7| = 1 \implies P(S=7) = \frac{1}{6} \approx 16.67\%$
\end{itemize}

\textbf{Conclusión combinatoria:} Aunque cada evento elemental tiene $P(x) = \frac{1}{6}$, el subconjunto de eventos cuya suma es $5$ tiene el \textbf{doble de densidad de probabilidad} ($33.33\%$) que cualquier otra suma, simplemente porque contiene más elementos de $\Omega$.

\medskip
\hrule
\medskip

\subsection*{3. ¿Qué es entonces la "Inferencia" a partir de Datos Históricos?}

Imaginemos que observamos un historial de $N = 100$ sorteos de esta micro-lotería:

$$D = \{x_1, x_2, \dots, x_{100}\}$$

La \textbf{inferencia estadística} consiste en realizar dos tareas formales con respecto a $\Omega$:

\subsubsection*{A. Verificación empírica del proceso generador}
Contrastar si la distribución empírica observada en los datos históricos $\hat{P}(S)$ converge a la distribución teórica $P(S)$ dictada por la combinatoria:
\begin{itemize}
    \item Si en el historial la suma 5 ocurre cerca del $33\%$ de las veces, confirmamos que el sistema responde a la física del espacio muestral sin sesgos mecánicos dominantes.
\end{itemize}

\subsubsection*{B. Reducción condicional del espacio de decisión}
Si tenemos que seleccionar un candidato para el evento $X_{t+1}$, la inferencia condiciona el espacio muestral a la región de \textbf{mayor densidad de masa probabilística}:

$$\Omega_{\text{óptimo}} = \{ x \in \Omega \mid x \in E_5 \} = \big\{ \{1,4\}, \{2,3\} \big\}$$

Cualquier elección fuera de $\Omega_{\text{óptimo}}$ pertenece a subconjuntos con menor probabilidad acumulada ($|E_s|/|\Omega| = 1/6$).

\subsubsection*{C. Inferencia Bayesiana / Ponderación Empírica}
Si en los datos observados $D$, la bola $4$ apareció en 45 de los 100 eventos y la bola $1$ solo en 10, podemos formular dos hipótesis:
\begin{enumerate}
    \item \textbf{Hipótesis nula ($H_0$):} El proceso es uniforme puro y la diferencia es varianza estocástica.
    \item \textbf{Hipótesis con sesgo/memoria empírica:} Ajustamos los pesos de la verosimilitud $P(D \mid \theta)$ para actualizar la distribución a priori de cada bola $\theta = (p_1, p_2, p_3, p_4)$, asignando mayores probabilidades de selección a combinaciones que involucren parámetros con mayor soporte de verosimilitud.
\end{enumerate}

\medskip
\hrule
\medskip

\subsection*{4. Traslación directa al caso real de Melate}

El principio matemático en tu aplicación es idéntico, solo cambia la escala de $\Omega$:

\begin{itemize}
    \item En Melate:
    $$|\Omega| = \binom{56}{6} = 32,468,436$$
    \item Cada combinación elemental $x$ tiene exactamente:
    $$P(x) = \frac{1}{32,468,436} \approx 3.08 \times 10^{-8}$$
    \item Sin embargo, si particionamos $\Omega$ mediante la variable aleatoria $S(x) = \sum_{i=1}^6 r_i$:
    \begin{itemize}
        \item La región $A = \{x \in \Omega \mid 125 \le S(x) \le 215\}$ concentra aproximadamente \textbf{$26$ millones de combinaciones}, representando el \textbf{$\approx 80\%$ de la masa probabilística} total.
        \item La región $B = \{x \in \Omega \mid S(x) < 80 \text{ ó } S(x) > 260\}$ contiene menos del \textbf{$1\%$} de las combinaciones posibles.
    \end{itemize}
\end{itemize}

\subsection*{Conclusión conceptual}

\textbf{Inferir} en tu aplicación no significa predecir el número que saldrá. Significa:
\begin{enumerate}
    \item Usar el historial $D$ para estimar parámetros de frecuencia y retardo empíricos.
    \item Restringir el espacio muestral $\Omega$ a los subconjuntos cuya cardinalidad $|A|$ maximiza $P(A) = \frac{|A|}{|\Omega|}$ bajo variables aleatorias agregadas (suma gaussiana y paridad).
    \item Muestrear candidatos exclusivamente dentro de ese subespacio de alta densidad.
\end{enumerate}

\end{document}